% source: 522c0ecf48be69427ae1d11937a9a679a69eeac1
\chapter{模 LWE（MLWE）}
\label{ch:mlwe}

\section{基于格的密码的三代}

我们已经看到了基于格的密码发展中的三个主要阶段。

\begin{itemize}
    \item \textbf{第一代：LWE}。其方程为
    \[
    b = A s + e,
    \]
    其中 $A$ 是一个随机矩阵。其安全性证明很优雅，但表示规模大，计算代价高。

    \item \textbf{第二代：RLWE}。矩阵被多项式环结构所取代。计算变得快得多，但代数结构变得更强、更特殊。

    \item \textbf{第三代：模 LWE}。这是规模大但简单的 LWE 构造与紧凑但高度结构化的 RLWE 构造之间的折中。
\end{itemize}

因此，历史路径为：

\begin{center}
\begin{tikzpicture}[node distance=9mm, every node/.style={font=\small}]
  \node[flownode] (a) {LWE};
  \node[flownode, right=of a] (b) {RLWE};
  \node[keynode, right=of b] (c) {Module-LWE};
  \draw[flowarrow] (a) -- (b);
  \draw[flowarrow] (b) -- (c);
\end{tikzpicture}
\end{center}

\section{什么是模？}

“模”（module）一词并不是随意的行话；它是一个精确的代数结构，而 Kyber 的秘密就位于其中。它推广了第~\ref{ch:linalg}~章的向量空间，允许标量来自环而不是域。

\begin{definition}[环上的模]
设 $R$ 是一个具有单位元 $1$ 的交换环。一个（左）\emph{$R$-模}是一个阿贝尔群 $(M,+)$，连同一个标量乘法 $R\times M\to M$（记作 $(r,m)\mapsto rm$），使得对所有 $r,s\in R$ 和 $x,y\in M$ 有：
\begin{enumerate}
  \item $r(x+y)=rx+ry$；
  \item $(r+s)x=rx+sx$；
  \item $(rs)x=r(sx)$；
  \item $1\,x=x$。
\end{enumerate}
\end{definition}

这些恰好就是把域换成环之后的向量空间公理。当 $R$ 是域时，$R$-模\emph{就是}向量空间；模是把“域”放宽为“环”所得到的严格更一般的对象。若干性质得以保留，而有一项便利之处丧失了。

\begin{itemize}
  \item \textbf{自由模与秩。}集合 $R^{k}=\{(m_1,\dots,m_k):m_i\in R\}$ 在逐分量运算下是一个 $R$-模，称为\emph{秩为 $k$ 的自由模}。与向量空间一样，它有一组基——标准单位向量 $e_1,\dots,e_k$——因此每个元素都能唯一地表示为它们的 $R$-组合。Kyber 的秘密空间恰好就是自由模 $R_q^{k}$。
  \item \textbf{矩阵作为同态作用。}矩阵 $A\in R^{k\times k}$ 定义了一个 $R$-线性映射 $R^{k}\to R^{k}$，$\mathbf{x}\mapsto A\mathbf{x}$，称为\emph{$R$-模同态}。模 LWE 映射 $\mathbf{s}\mapsto A\mathbf{s}$ 就是其中之一。
  \item \textbf{子模与理想。}子模是子空间的类比；把 $R$ 看作自身上的模时，其子模恰好就是它的\emph{理想}（第~\ref{ch:polyring}~章）——这就是为什么说 RLWE 建立在\emph{理想}格之上。
  \item \textbf{丧失了什么。}在一般的环上，并非每个模都是自由的，也并非每个子模都有基——这是与向量空间的一个实质性区别。Kyber 通过在自由模 $R_q^{k}$ 中工作而完全避开了这一点。
\end{itemize}

精确地命名了这一结构之后，三代演进就成了关于\emph{秘密位于哪个模中}的陈述：
\[
\underbrace{s\in\mathbb{Z}_q^{n}}_{\text{LWE}}
\qquad
\underbrace{s(x)\in R_q}_{\text{RLWE}}
\qquad
\underbrace{\mathbf{s}\in R_q^{k}}_{\text{Module-LWE}}.
\]
LWE 使用域元素组成的向量；RLWE 使用单个环元素；模 LWE 使用环 $R_q$ 上长度为 $k$ 的向量——一个秩为 $k$ 的自由模。

\section{Kyber 的秘密}

在 ML-KEM（Kyber）中，秘密不是单个多项式，而是一个多项式向量。安全级别决定了该向量的维数。

\begin{itemize}
    \item Kyber512 使用 $k=2$。
    \item Kyber768 使用 $k=3$。
    \item Kyber1024 使用 $k=4$。
\end{itemize}

因此，安全级别与模的大小相关联。

\section{为什么矩阵又回来了？}

乍看之下，人们会惊讶于在 Kyber 中再次看到矩阵。但这个矩阵并不是整数上的大型随机矩阵，其元素是多项式。

例如，对于 $k=2$，我们可以有
\[
A =
\begin{bmatrix}
 a_{11}(x) & a_{12}(x) \\
 a_{21}(x) & a_{22}(x)
\end{bmatrix}.
\]

每个元素都是环 $R_q$ 中的一个多项式。这个对象仍然是矩阵，但矩阵元素现在是环元素，而不是普通的数。

\section{Kyber 公钥生成}

公钥生成在思想上与 LWE 非常相似。我们选取一个秘密向量 $\mathbf{s}$、一个噪声向量 $\mathbf{e}$ 以及一个多项式矩阵 $A$，然后计算
\[
\mathbf{t} = A\mathbf{s} + \mathbf{e}.
\]

这与 LWE 的支配思想完全相同，只是标量现在变成了环元素。下图针对最小情形 $k=2$ 具体展示了 $A$ 与 $\mathbf{s}$ 的形状（下标从 $0$ 到 $k-1$，与下面的 Sage 代码一致）。

\begin{figure}[H]
\centering
\begin{tikzpicture}[every node/.style={font=\small}]
  % A 2x2 matrix of ring elements: each entry a_{ij}(x) is itself a
  % polynomial in R_q, not a single scalar as it would be for plain LWE
  \node[draw=pqcblue, thick, fill=pqclight, minimum width=2.0cm, minimum height=1.0cm] (a00) at (0,1.1) {$a_{00}(x)$};
  \node[draw=pqcblue, thick, fill=pqclight, minimum width=2.0cm, minimum height=1.0cm] (a01) at (2.2,1.1) {$a_{01}(x)$};
  \node[draw=pqcblue, thick, fill=pqclight, minimum width=2.0cm, minimum height=1.0cm] (a10) at (0,-0.1) {$a_{10}(x)$};
  \node[draw=pqcblue, thick, fill=pqclight, minimum width=2.0cm, minimum height=1.0cm] (a11) at (2.2,-0.1) {$a_{11}(x)$};
  \node[fit=(a00)(a01)(a10)(a11), inner sep=6pt, draw=pqcgray, thick] (Abox) {};
  \node[pqcblue, font=\small] at ($(Abox.south)+(0,-0.5)$) {$A \in R_q^{2\times 2}$};

  % A length-2 vector of ring elements: the secret s, one entry per row of A
  \node[draw=pqcorange, thick, fill=white, minimum width=2.0cm, minimum height=1.0cm] (s0) at (5.6,1.1) {$s_0(x)$};
  \node[draw=pqcorange, thick, fill=white, minimum width=2.0cm, minimum height=1.0cm] (s1) at (5.6,-0.1) {$s_1(x)$};
  \node[fit=(s0)(s1), inner sep=6pt, draw=pqcgray, thick] (sbox) {};
  \node[pqcorange, font=\small] at ($(sbox.south)+(0,-0.5)$) {$\mathbf{s} \in R_q^{2}$};
\end{tikzpicture}
\caption{$k=2$ 时模结构的具体形态：公开矩阵 $A \in R_q^{2\times2}$（蓝色，左）有四个元素 $a_{00}(x),a_{01}(x),a_{10}(x),a_{11}(x)$，每个都是 $R_q$ 中的完整多项式；秘密 $\mathbf{s}=(s_0(x),s_1(x))$（橙色，右）是由这样的多项式组成的长度为 $2$ 的向量。与前面的“秘密对象”一栏对照：普通 LWE 会在每个方框中放一个标量，而 RLWE 会把整幅图压缩成仅一个方框。}
\label{fig:mlwe-module-structure}
\end{figure}

\section{为什么这比普通 LWE 更好？}

我们可以比较这三种设计。

\begin{itemize}
    \item \textbf{LWE}：矩阵巨大，公钥很大。
    \item \textbf{RLWE}：表示紧凑，但代数结构非常强。
    \item \textbf{模 LWE}：矩阵小，元素是多项式，在紧凑性与灵活性之间取得了非常好的折中。
\end{itemize}

换言之，模 LWE 使用一个小的多项式矩阵，而不是一个巨大的整数矩阵或单个大多项式。

\section{用 Sage 构造一个玩具版模}

我们可以在 Sage 中构造一个非常小的玩具版示例。首先定义环。

\begin{lstlisting}[language=Python]
q = 17
R.<x> = PolynomialRing(GF(q))
S = R.quotient(x^8 + 1)
\end{lstlisting}

现在把秘密定义为由两个多项式组成的向量。

\begin{lstlisting}[language=Python]
from random import choice
def small_poly():
    return S([choice([-1, 0, 1]) for _ in range(8)])

s = vector(S, [small_poly(), small_poly()])
print(s)
\end{lstlisting}

构造一个小的多项式矩阵。

\begin{lstlisting}[language=Python]
A = Matrix(S, [[S.random_element(), S.random_element()],
               [S.random_element(), S.random_element()]])
print(A)
\end{lstlisting}

构造一个噪声向量。

\begin{lstlisting}[language=Python]
e = vector(S, [small_poly(), small_poly()])
assert all(c in [0, 1, 16] for p in list(s) + list(e) for c in p.lift().list())
\end{lstlisting}

然后计算类似公钥的值。

\begin{lstlisting}[language=Python]
t = A * s + e
print(t)
\end{lstlisting}

这具有模 LWE 样本的代数形状。公开矩阵是均匀分布的，而秘密和误差使用小系数；它仍然只是一个玩具版示例，而不是 FIPS 203 的采样器。

\section{模 LWE 问题}

有了模结构，模 LWE 就是把第~\ref{ch:lwe}~章的 LWE 分布中的标量从 $\mathbb{Z}_q$ 提升到环 $R_q$ 而得到的分布。

\begin{definition}[模 LWE]
固定环 $R_q=\mathbb{Z}_q[x]/(x^n+1)$、模秩 $k$，以及 $R_q$ 上集中于小元素的误差分布 $\chi$。对于秘密 $\mathbf{s}\in R_q^{k}$，\emph{模 LWE 分布}输出
\[
(\mathbf{a},\,b),\qquad \mathbf{a}\leftarrow R_q^{k}\ \text{uniform},\quad e\leftarrow\chi,\quad b=\langle \mathbf{a},\mathbf{s}\rangle + e\in R_q.
\]
把 $m$ 个样本作为 $A\in R_q^{m\times k}$ 的行堆叠起来，便得到矩阵形式 $(A,\ \mathbf{t}=A\mathbf{s}+\mathbf{e})$。\emph{搜索 MLWE} 是要恢复 $\mathbf{s}$；\emph{判定 MLWE} 是要区分 $(A,\,A\mathbf{s}+\mathbf{e})$ 与均匀的二元组 $(A,\,\mathbf{u})$。
\end{definition}

有两个特殊情形界定了这一定义的两端：$k=1$ 即为 RLWE（单个环元素），而 $n=1$ 且 $R_q=\mathbb{Z}_q$ 则还原为普通 LWE。模 LWE 在二者之间插值，并继承了与 LWE 同类的最坏情形困难性保证：在平均情形下求解它至少与\emph{模}格上的标准问题一样困难~\cite{langlois2015}。在 Kyber 中矩阵是方阵（$m=k$），从而得到紧凑形式
\[
\mathbf{t} = A\mathbf{s} + \mathbf{e}, \qquad A\in R_q^{k\times k},\ \ \mathbf{s},\mathbf{e}\in R_q^{k}.
\]

\section{为什么模 LWE 被视为一种折中}

许多人误以为模 LWE 之所以“更好”，只是因为它使用了更多的结构。这不是正确的理解。真正的要点在于，模 LWE 削弱了纯 RLWE 最强的代数结构，同时保留了其效率优势。换言之，它在以下几者之间进行权衡：

\begin{itemize}
    \item 普通 LWE 的巨大规模，
    \item RLWE 的强结构，
    \item 以及 MLWE 所提供的实用中间地带。
\end{itemize}

这种平衡是 Kyber 设计的核心。

\section{为什么 Kyber 的参数是 $k=2,3,4$}

取值 $k=2,3,4$ 对应模的不同维数，恰好就是 Kyber 的三个安全级别~\cite{kyber2018,langlois2015}。在实践中：

\begin{itemize}
    \item $k=2$ 给出 Kyber512，
    \item $k=3$ 给出 Kyber768，
    \item $k=4$ 给出 Kyber1024。
\end{itemize}

增大 $k$ 会增大模的规模，从而提高安全级别。环的维数和模数保持不变，但标准化的 ML-KEM 参数集也使用了不同的参数元组：特别地，对于 ML-KEM-512、-768 和 -1024，$\eta_1$ 分别为 $3,2,2$，$(d_u,d_v)$ 分别为 $(10,4),(10,4),(11,5)$~\cite{fips203}。

\section{一个完整的玩具版 MLWE 实验}

一个有益的练习是完整地实现一个玩具版 MLWE。

\begin{enumerate}
    \item 定义环 $R_q$。
    \item 构造一个小的多项式矩阵 $A$。
    \item 构造一个多项式向量 $\mathbf{s}$。
    \item 构造一个噪声向量 $\mathbf{e}$。
    \item 计算 $\mathbf{t}=A\mathbf{s}+\mathbf{e}$。
    \item 打印结果并考察其结构。
\end{enumerate}

这能让人切实理解为什么 Kyber 可以看作 LWE 的一种结构化推广。用伪代码写出，并沿用上面 Sage 代码中的名称 $A$、$\mathbf{s}$、$\mathbf{e}$、$\mathbf{t}$，同样的六个步骤即为算法~\ref{alg:mlwe-sample}。

\begin{algorithm}[H]
\caption{模 LWE 样本生成}
\label{alg:mlwe-sample}
\begin{algorithmic}[1]
\Require 环 $R_q = \mathbb{Z}_q[x]/(x^n+1)$；模维数 $k$
\Ensure MLWE 样本 $(A,\mathbf{t})$；秘密 $\mathbf{s}$
\State $A \gets$ $R_q$ 上均匀随机的 $k\times k$ 矩阵
\State $\mathbf{s} \gets$ 由 $R_q$ 中小（噪声量级）元素组成的长度为 $k$ 的向量 \Comment{秘密}
\State $\mathbf{e} \gets$ 由 $R_q$ 中小（噪声量级）元素组成的长度为 $k$ 的向量 \Comment{误差}
\State $\mathbf{t} \gets A\mathbf{s} + \mathbf{e}$ \Comment{由 $k$ 个环元素组成的向量，而非单个标量}
\State \Return 以 $(A,\mathbf{t})$ 作为样本，$\mathbf{s}$ 作为秘密
\end{algorithmic}
\end{algorithm}

实际的 ML-KEM 具有这种算法形状，其中 $n=256$、$k=2,3,4$，但各标准化参数集之间的差异不止在于 $k$；它们的采样参数和压缩参数在 FIPS~203 中规定。

\section{小结}

关键的教训是：模 LWE 并不只是“带有更多多项式的 LWE”。它是两个极端之间的中间地带：

\begin{table}[h]
\centering
\begin{tabular}{|c|c|c|}
\hline
方案 & 秘密对象 & 主要思想 \\
\hline
LWE & $\mathbb{Z}_q$ 上的向量 & 简单通用，但规模大 \\
\hline
RLWE & $R_q$ 上的单个多项式 & 紧凑快速，但高度结构化 \\
\hline
模 LWE & $R_q$ 上的多项式向量 & 均衡的折中，为 Kyber 所采用 \\
\hline
\end{tabular}
\end{table}

重要的洞见在于，Kyber 并不只是换了个名字的 RLWE。它是同一带噪声线性结构的模版本，其组织方式在保持强安全假设的同时带来了实用的效率。

\printbibliography[heading=chapterrefs]
